% Einheit 4 von 4 – Lot und Extremum (bank/orthogonalitaet/e4.jsonl, zusammenbau v0.1)
%% TODO Zeitmarke Einheit 4: Marken-Zeile nicht lesbar
%% TODO Zweigzeile Teil 1: was hier gelernt wird (Satz in Schülersprache aus dem Einheitstitel) – Einheit 4
\einheitenkopf[e4]{Einheit 4 von 4 · Lot und Extremum}
\zweigzeile{keine Prüfungsaufgabe}
\setcounter{aufgabe}{17}

%% TODO Titel: Ich-kann-Satz fehlt in der Bank (Platzhalter: Kettenname) – Nr. 18
%% TODO Anweisung: ein Satz über der ersten Teilaufgabe fehlt in der Bank – Nr. 18
\begin{aufgabe}{Lot und Extremum}
%% TODO \rechenplatz{n}: Schrittzahl des Grundfalls fehlt in der Bank (nur bei fester Schreibform, 2.3 b)
\begin{teile}
\teil Von $P$ wird das Lot auf die Gerade $g$ gefällt, $F$ ist der Lotfußpunkt. Welches Paar ist senkrecht? \\ \kreuz{$\overrightarrow{PF}$ und der Richtungsvektor von $g$} \\ \kreuz{$\overrightarrow{PF}$ und der Ortsvektor von $F$} \\ \kreuz{der Ortsvektor von $P$ und der Richtungsvektor von $g$}
\teil Gegeben sind $P(3 | 0 | 2)$ und $g\colon \vec x = (0 | 1 | 0) + \lambda \cdot (2 | 1 | 2)$. Bestimme den Punkt $F$ von $g$, der $P$ am nächsten liegt, und den Abstand von $P$ zu $g$. Begründe deinen Ansatz.
\teil Gegeben sind $C(-3 | 0 | 2)$, $D(3 | 0 | 2)$ und für $0 < k \le 6$ die Punkte $F_k(0 | k | 12 - 2k)$; alle Dreiecke $CDF_k$ sind gleichschenklig mit der Basis $CD$. Bestimme den Wert von $k$, für den der Innenwinkel bei $F_k$ am größten ist, und erläutere deinen Weg \hfill (Abitur 2025 LK).
\teil Gegeben ist ein Quadrat $ABCD$. Die Gerade $g$ verläuft durch $A$ und den Mittelpunkt $M$ der Seite $BC$; $F$ ist der Lotfußpunkt von $B$ auf $g$. Begründe ohne Koordinaten, dass $|AF| = 2 \cdot |BF|$ gilt \hfill (Abitur 2022 LK).
\end{teile}
\end{aufgabe}

%% TODO Titel: Ich-kann-Satz fehlt in der Bank (Platzhalter: Kettenname) – Nr. 19
%% TODO Anweisung: ein Satz über der ersten Teilaufgabe fehlt in der Bank – Nr. 19
\begin{aufgabe}{Flächengleichung des flächenkleinsten Dreiecks begründen}
\begin{teile}
\teil Gegeben sind $R(0 | 1 | 3)$, $S(0 | 5 | 3)$ und die Gerade $t\colon \vec x = (1 | 3 | 0) + \lambda \cdot (2 | 0 | 1)$. Jeder Punkt $T_\lambda$ von $t$ hat von $R$ und $S$ denselben Abstand. Unter den Dreiecken $RST_\lambda$ hat eines den kleinsten Flächeninhalt. Begründe, dass der zugehörige Wert von $\lambda$ die Gleichung $(1 + 2\lambda | 0 | \lambda - 3) \circ (2 | 0 | 1) = 0$ löst \hfill (Abitur 2020 GK).
\end{teile}
\end{aufgabe}

%% TODO Titel: Ich-kann-Satz fehlt in der Bank (Platzhalter: Kettenname) – Nr. 20
%% TODO Anweisung: ein Satz über der ersten Teilaufgabe fehlt in der Bank – Nr. 20
\begin{aufgabe}{Lot und Extremum – Fehler finden, Begründen, Anwendung}
\begin{teile}
\teil Für die Dreiecke $RST_\lambda$ mit fester Basis $RS$ und Spitze $T_\lambda$ auf einer Geraden $t$ sucht Kim das flächenkleinste und setzt an: „Die Fläche ist am kleinsten, wenn $|RT_\lambda|$ am kleinsten ist.“ Finde den Fehler.
\teil Begründe, warum unter allen Strecken von einem Punkt $P$ zu einer Geraden $g$ die senkrechte die kürzeste ist.
\teil Ein Haus steht bei $H(2 | 5 | 0)$, eine gerade Straße verläuft entlang $s\colon \vec x = (0 | 1 | 0) + \lambda \cdot (2 | 1 | 0)$ (1 LE = 100 m). Wo soll der Zufahrtsweg auf die Straße treffen, damit er möglichst kurz ist, und wie lang wird er?
\end{teile}
\end{aufgabe}
